\documentclass[../main.tex]{subfiles}
\begin{document}
% \section{Utilisation  \tkzname{pgfmath} et de \tkzname{fp.sty} }
\section{\tkzname{pgfmath}和\tkzname{fp.sty}工具}
%--------------------------------------------------------------------------->
% \subsection{\tkzname{pgfmath}}
\subsection{\tkzname{pgfmath}}

% On peut faire maintenant beaucoup de tracés sans gnuplot, voici à titre d'exemple et d'après une idée d'Herbert Voss (le membre le plus actif de la communauté Pstricks) un exemple de courbes obtenues avec seulement Tikz.
即便是没有Gnuplot，也可以绘制图形，下面是\tkzname{Herbert Voss}(Pstricks社区最活跃的成员)仅使用\TIKZ{}绘图的一种方法。

\begin{center}
\begin{tkzexample}[]
\begin{tikzpicture}
  \def\Asmall{0.7 } \def\Abig{3 } \def\B{20}%Herbert Voss
  \path[fill=blue!40!black,domain=-pi:pi,samples=500,smooth,variable=\t]%
          plot({\Abig*cos(\t r)+\Asmall*cos(\B*\t r)},%
               {0.5*\Abig*sin(\t r)+0.5*\Asmall*sin(\B*\t r)});
  \def\Asmall{0.7 } \def\Abig{3 } \def\B{10}
  \path[shift={(1,1)},fill=blue!40!black,%
        domain=-pi:pi,samples=500,smooth,variable=\t]%
        plot({\Abig*cos(\t r)+\Asmall*cos(\B*\t r)},%
             {0.5*\Abig*sin(\t r)+0.5*\Asmall*sin(\B*\t r)});
\end{tikzpicture}
\end{tkzexample}
\end{center}

% \subsection{\tkzname{fp.sty}}
\subsection{\tkzname{fp.sty}}

% Le principal problème  de \tkzname{fp.sty} se produit lors de l'évaluation par exemple de $(-4)^2$ ce qui peut se traduire avec fp par~:
\tkzname{fp.sty}的主要问题是可能会造成计算错误，如：$(-4)^2$的计算，
% 在代码中，可以先计算出该值，再在后续代码中使用该计算结果。

\begin{tkzltxexample}[]
\begin{tikzpicture}
   \FPeval\result{(-4)^2}
\end{tikzpicture}
\end{tkzltxexample}

% ce qui donne une erreur car fp utilise les logarithmes pour faire cette évaluation. \tkzname{tkz-fct.sty} modifie la macro \tkzcname{FP@pow} pour éviter cette erreur
由于使用对数的方式进行计算，因此，这可能会造成错误。
为此，\tkzname{tkz-fct.sty}宏包通过修改\tkzname{FP@pow}宏以避免这个错误。

%  Pour calculer les pentes des tangentes et pour placer des points sur les courbes, mon module traduit l'expression donnée pour gnuplot et la stocke dans une commande \tkzcname{tkzFcta}, pour être utilisée ensuite avec les macros \tkzcname{tkzDefPointByFct}\ et \tkzcname{tkzDrawTangentLine}.
当需要计算切线的斜率时并在函数曲线上放置点时，该宏包通过转换表达式为\tkzname{gnuplot}提供数据，
并将结果保存在\tkzname{tkzFcta}宏中，
在后续代码中，可以使用\tkzcname{tkzDefPointByFct} 和 \tkzcname{tkzDrawTangentLine}命令使用这些数据。

% mais si vous voulez placer un point de ce graphe ayant pour abscisse $x=2$, il est alors préférable de choisir la première méthode.
不过，如果需要在该图中放置一个$x=2$的点，建议使用第一种方法比较合理。

% Sinon pour une fonction polynômiale, il sera nécessaire pour utiliser les macros relatives aux images et aux tangentes de mettre le polynôme sous la forme d'Horner.
% Ainsi avec \tkzcname{tkzFct}, l'argument $x^4-2x^3+4x-5$ peut être écrit : |-5+x*(0.5+4*x*(x*(-2+x*1)))|.
否则，对于多项式函数，在使用与图像切线相关的命令时，需要将多项式表示为\tkzname{Horner}的形式。
因此，在使用\tkzcname{tkzFct}命令时，需要将$x^4-2x^3+4x-5$改写成：
|-5+x*(0.5+4*x*(x*(-2+x*1)))|.\

% Voici ce qu'il faut donc faire :
因此，需要做的是：

\begin{center}
\begin{tkzexample}[]
\begin{tikzpicture}
 \tkzInit[xmin=-0.2,xmax=0.2,xstep=.1,
          ymin=-12,ymax=6,ystep=2]
 \tkzGrid
 \tkzAxeXY
 \tkzFct[domain = -.1:.2]%
 {-5+x*(0.5+4*x*(x*(-2+x*1)))}
\end{tikzpicture}
\end{tkzexample}
\end{center}

\end{document}
\endinput
